\section{An explicit arithmetic fiber}\label{sec:example}

Consider the squarefree quintic
\[
 P_5(T)=(T^3-19)(T^2+T+1)
 =T^5+T^4+T^3-19T^2-19T-19.
\]
Berend and Bilu introduced this polynomial in their study of polynomials
having a root modulo every integer \cite{BerendBilu}.
It has no rational root: the rational-root theorem excludes a root of
\(T^3-19\), while \(T^2+T+1\) has discriminant \(-3\).  On the other hand,
\(\Spec\Q[T]/(P_5)\) has a point over every completion of \(\Q\).  The cubic
factor has a real root.  For a prime
\(p\notin\{2,3,19\}\), if \(p\equiv2\pmod3\), cubing is an automorphism of
\(\mathbb F_p^\times\), so \(T^3-19\) has a simple root modulo \(p\); if
\(p\equiv1\pmod3\), then \(\mathbb F_p^\times\) contains a nontrivial cube
root of unity, giving a simple root of \(T^2+T+1\).  At the exceptional
primes, Hensel witnesses are
\[
\begin{array}{c|c|c}
p&\text{factor and approximate root}&\text{check}\\ \hline
2&T^3-19,\ T=1&-18\equiv0\pmod2,\quad 3\in\mathbb Z_2^*\\
3&T^3-19,\ T=-2&v_3(-27)=3>2v_3(12)=2\\
19&T^2+T+1,\ T=7&57\equiv0\pmod{19},\quad15\in\mathbb Z_{19}^*.
\end{array}
\]

We now realize this algebra by the formulas above.  The translation \(a=0\)
is admissible, and
\[
 G(S)=P_5(S)-P_5(0)
 =S^5+S^4+S^3-19S^2-19S.
\]
Thus \((g_1,g_2,g_3,g_4,g_5)=(-19,-19,1,1,1)\).  Put
\[
 t=1+xy,\qquad q=t^2z-19y^2(1+3t),
\]
and define \(F=(F_1,F_2,F_3)\) by
\begin{align}
 F_1={}&tq,\notag\\
 F_2={}&19y-3xq+38tq-4t^2x^2q^4-5t^2x^3q^5,\notag\\
 F_3={}&19x(5-3t)+x^3z+2(xq)^4+3(xq)^5.
 \label{eq:explicit-quintic-map}
\end{align}
This is an output rescaling of the quadratic-gauge map, so
\[
 \det DF=-722,\qquad \gdeg(F)=5.
\]
At the target \(y=(1,0,-38)\), the inverse equation is
\[
 G(S)-\frac{-19}{2}(-2)=G(S)-19=P_5(S).
\]
The squarefree fiber theorem therefore gives
\[
 \boxed{F^{-1}(1,0,-38)
 \simeq\Spec\Q[T]/\bigl((T^3-19)(T^2+T+1)\bigr).}
\]
This full Keller fiber has points over \(\mathbb R\) and every
\(\mathbb Q_p\), but no rational point.  The displayed map, literal fiber,
rank, finite \'etaleness, and these local and global point statements are
also included in the Lean development.  Composing \(F\) with the diagonal
target scaling
\((u,v,w)\mapsto(u,-v/38,w/19)\) gives determinant one and sends the target
back to \((1,0,-2)\).
